% Part: proof-theory
% Chapter: proof-search
% Section: introduction

\documentclass[../../../include/open-logic-section]{subfiles}

\begin{document}

\olfileid{pt}{ps}{int}

\olsection{مقدمة}

تحدد مسلّمات وقواعد استدلال أنظمة~!!{proof} مثل~\Log{G3c} أو~\Log{N1i}
أي أشجار المتتاليات أو~!!{formula}s تُعد~!!{proof}s صحيحة. فإذا أُعطيت شجرة
كهذه من المتتاليات أو~!!{formula}s (وربما معها معلومات إضافية، مثل القواعد
المطبقة في كل خطوة، أو وسوم الفروض وقواعد الاستدلال التي تفرّغها)، أتاحت لنا
تعريفات المسلّمات والقواعد التحقق مما إذا كانت الشجرة المعطاة~!!{proof}
صحيحًا. أما \emph{العثور} على~!!a{proof} لمتتالية معطاة، أو لـ~!!{formula}
معطاة من فروض معينة، فمسألة مختلفة وأصعب عادةً. وهذه هي مسألة \emph{البحث
عن البرهان}.

توجد طريقة ساذجة جدًا للبحث عن البرهان. يمكن عد~!!{formula}s متتاليات من
الرموز المأخوذة من أبجدية منتهية. ويمكن كذلك عد المتتاليات وأشجار
المتتاليات أو~!!{formula}s متتاليات من~!!{formula}s (مع أقواس خاصة، ربما،
للدلالة على تشعب الشجرة). وتتيح لنا المسلّمات والقواعد فحص ما إذا كانت
متتالية رموز تمثل~!!{proof} صحيحًا فحصًا آليًا. ولذلك نستطيع توليد
\emph{جميع}~!!{proof}s الصحيحة الممكنة آليًا (بتوليد جميع متتاليات الرموز
من الأبجدية المستعملة لتمثيل~!!{proof}s، ثم فحص كل متتالية مرشحة لمعرفة ما
إذا كانت تمثل~!!{proof} صحيحًا). وإجراء البحث الساذج هو أن نولّد جميع
!!{proof}s الصحيحة إلى أن نجد واحدًا يكون~!!{proof} للمتتالية المعطاة (أو
لـ~!!{formula} المعطاة، مع فروض مفتوحة ضمن الفروض المعطاة).

والطريقة الأفضل هي استعمال الأسلوب الساذج الذي تتبعه يدويًا للعثور على
!!a{proof}: تبدأ بالمتتالية النهائية، وتختار~!!a{formula}، وتطبق قاعدة
استدلال ``إلى الخلف'' لتولّد مقدمة أو مقدمات لو كان لها~!!{proof}s لأمكن
جمعها مع الاستدلال الأخير في~!!a{proof} للمتتالية المعطاة. وستتبع طريقة
مشابهة، وإن كانت أعقد، للعثور على~!!a{proof} في نظام استنباط طبيعي.

لكن هذه الطريقة ليست مضمونة: فقد تصل إلى طريق مسدود إذا اخترت القاعدة
الخاطئة أو رتبت~!!{formula}s ترتيبًا خاطئًا. كما أنها غير محددة تمامًا، إذ
قد يوجد في كل نقطة أكثر من~!!{formula} واحدة يمكن اختيارها، أو أكثر من
قاعدة يمكن تطبيقها إلى الخلف. و\emph{خوارزمية البحث عن البرهان} إجراء محدد
تمامًا يجد~!!a{proof} متى وُجد (أي إنه مضمون).

تلائم بعض أنظمة~!!{proof} خوارزميات بحث أبسط من غيرها. ففي أنظمة
المتتاليات، يحدد~!!{main operator} لـ~!!a{formula} والجانب الذي تقع فيه
القاعدة التي ينبغي تطبيقها إلى الخلف. فمثلًا، إذا أردت العثور على
!!a{proof} للمتتالية~$\Gamma \Sequent \Delta,!A\land !B$ التي تقع فيها
$!A\land !B$ على اليمين، فعليك تطبيق~\RightR{\land} إلى الخلف وتوليد
المقدمتين المقابلتين~$\Gamma \Sequent \Delta,!A$ و
$\Gamma \Sequent \Delta,!B$. أما إذا كان نظامك~\Log{G1c}، فإن وجود
الانقباض يعقّد الأمر لأنه يمنحك خيارات أكثر: يمكنك أيضًا تطبيق
\RightR{\Contraction} إلى الخلف وتوليد المقدمة
$\Gamma \Sequent \Delta,!A\land !B,!A\land !B$، ثم المقدمة ذات النسخ
الثلاث، وهكذا. ويزداد الأمر سوءًا في~\Log{G2c}. فتطبيق~\RightR{\land} إلى
الخلف يقتضي أيضًا تخمين~$\Gamma_1,\Gamma_2$ بحيث
$\Gamma=\Gamma_1\cup\Gamma_2$، وتخمين~$\Delta_1,\Delta_2$ بحيث
$\Delta=\Delta_1\cup\Delta_2$، ثم تجربة المقدمتين
$\Gamma_1\Sequent\Delta_1,!A$ و$\Gamma_2\Sequent\Delta_2,!B$. وقد يقودك
تخمين خاطئ إلى طريق مسدود لاحقًا، فتضطر إلى الرجوع للبداية واختيار سياقات
مختلفة. وإذا كانت~!!{formula} هي~$!A\lor !B$، فعليك اختيار إحدى صورتي
\RightR{\lor}، فتنتج إحدى المقدمتين
$\Gamma\Sequent\Delta,!A$ أو~$\Gamma\Sequent\Delta,!B$. ويضطرك الاختيار
الخاطئ مرة أخرى إلى الرجوع.

لا يعاني~\Log{G3c} هذه المشكلات. فليس فيه قواعد بنيوية، ويتحدد اختيار قاعدة
الاستدلال بـ~!!{main operator} وبالجانب الذي تقع فيه~!!{formula}. ولذلك
يلائم~\Log{G3c} البحث عن البرهان أكثر من~\Log{G1c} أو~\Log{G2c}. وينتج عن
بحث ناجح~!!a{proof} في~\Log{G3c} يمكن بعد ذلك ترجمة~!!{proof} إلى
\Log{G1c} أو~\Log{G2c} (أو حتى~\Log{N1c}). ومن ثم فإن وجود خوارزمية بحث
لـ~\Log{G3c} وخوارزميات لترجمة~!!{proof}s~\Log{G3c} إلى~!!{proof}s في
الأنظمة الأخرى يحل مسألة البحث في تلك الأنظمة أيضًا.

كيف نعرف أن خوارزمية البحث مضمونة؟ يقتضي ذلك إظهار أنه إذا أخفقت الخوارزمية
في العثور على~!!a{proof}، استحال وجود~!!{proof}. فقد تخفق خوارزمية سيئة في
العثور على برهان مع أنه موجود. ولا تواجه الخوارزمية الساذجة هذه المشكلة
لأنها تبحث في جميع~!!{proof}s الممكنة، لكن المفروض في خوارزمية البحث أن
تكون فعّالة وأن تبذل أقل قدر من الجهد.

الفكرة الأساسية لإثبات ضمان خوارزمية البحث هي أن نستخرج من طريقة إخفاقها
دليلًا على استحالة وجود~!!a{proof} للمتتالية. ونفعل ذلك بإظهار أنها غير
صالحة. ففي منطق الرتبة الأولى الكلاسيكي، المتتاليات الصالحة
$\Gamma\Sequent\Delta$ هي التي يجعل فيها كل~!!{structure} !!{formula}
واحدةً على الأقل في~$\Gamma$ كاذبة، أو~!!{formula} واحدةً على الأقل في
$\Delta$ صادقة. والنسق~\Log{G3c} \emph{سليم}، أي إنه لا~!!{prove}s إلا
متتاليات صالحة. ويُثبت ذلك بالاستقراء في~!!{proof}s: فالمسلّمات صالحة
بوضوح، وإذا كانت مقدمات قاعدة صالحة كانت نتيجتها صالحة. ولإثبات ضمان
خوارزمية البحث، نبيّن أنه إذا أخفق البحث عن~!!a{proof} للمتتالية
$\Gamma\Sequent\Delta$، أمكن تعريف~!!a{structure} تكون فيه جميع
!!{formula}s في~$\Gamma$ صادقة وجميع~!!{formula}s في~$\Delta$ كاذبة. وهذا
يثبت أيضًا أن~\Log{G3c} \emph{كامل}؛ أي إنه إذا كانت
$\Gamma\Sequent\Delta$ صالحة، وُجد لها~!!{proof} في~\Log{G3c}. ولأن كل
بحث فاشل يبيّن أن المتتالية غير صالحة، فلا بد أن ينجح كل بحث عن برهان
لمتتالية صالحة.

\end{document}
